\chapter*{Mode d'emploi}
\addcontentsline{toc}{chapter}{Mode d'emploi}

Chaque chapitre suit une architecture stable : réactivation des prérequis,
automatismes courts, activité de recherche, cours structuré, exemples résolus,
exercices progressifs, problèmes de transfert, défis et bilan. Cette régularité
permet à l'élève de savoir immédiatement ce que l'on attend de lui.

\begin{multicols}{2}
\textbf{Niveaux de difficulté}
\begin{itemize}
  \item \diff{1} application directe et maîtrise attendue;
  \item \diff{2} combinaison de méthodes et raisonnement;
  \item \diff{3} recherche, preuve ou problème olympique.
\end{itemize}

\columnbreak
\textbf{Statut des contenus}
\begin{itemize}
  \item sans mention : parcours principal;
  \item encadré violet : pour aller plus loin;
  \item badge Python : activité numérique, jamais indispensable à la lecture.
\end{itemize}
\end{multicols}

Les automatismes de fin d'ouvrage sont volontairement mixtes : ils reprennent
des acquis plus anciens au lieu de ne travailler que la notion du jour. Pour
apprendre durablement, il faut retrouver une méthode après un délai et choisir
la bonne procédure parmi plusieurs possibilités.
